\DocumentMetadata{tagging=on,pdfstandard={UA-2,A-4f}}
\PassOptionsToPackage{colorlinks,linktocpage}{hyperref}
\documentclass{l3in2edoc}

\usepackage{fontspec}
\usepackage{lum-pagella}
\setmainfont{TeX Gyre Pagella}[Numbers=OldStyle]

\usepackage[it,medium]{titlesec}

\usepackage{array,vertbars}
\usepackage[svgnames]{xcolor}

\usepackage{geometry}
\geometry{b5paper}

\DeclareRobustCommand\pkg{\textsf}
\def\pkgopt#1{\texttt{[#1]}}

\usepackage{microtype}
\linespread{1.1}
\frenchspacing

\GetFileInfo{vertbars.sty}
\begin{document}

\title{The \pkg{vertbars} package}
\author{%
 Author: Peter Wilson, Herries Press\\
 Maintainer: \LaTeX{} Project\\
 \normalsize\url{https://github.com/LaTeX-Package-Repositories/herries-press}%
}
\date{\fileversion \qquad \filedate}

\maketitle

\section{Documentation}

Because this package uses \pkg{lineno} as a driver,
please read the documentation of that package first to
understand the mechanisms used here.

The \pkg{vertbars} package takes the same options as the \pkg{lineno} package.
In particular, the \texttt{switch} and \texttt{switch*} options control which side of the page any bars will be printed.
The package automatically calls the lineno package, so you just need to write:
\begin{verbatim}
    \usepackage[...]{vertbars}
\end{verbatim}

The package provides the \texttt{vertbar} environment, which is equivalent to the \texttt{linenumbers} environment except that a vertical bar replaces the line number.
Text paragraphs within a \texttt{vertbar} environment will be marked with a vertical bar in the margin.
Nested environments generate multiple marginal bars.

The width of the bars is controlled by the value of \cmd\barwidth, which is initialised to \verb|0.4pt|, and can be changed via \cmd\setlength.

The horizontal separation between adjacent bars is controlled by the value of \cmd\barspace, which is initialised to \verb|0.5\linenumbersep|, and can be changed via \cmd\setlength.
(\cmd\linenumbersep is a command from the \pkg{lineno} package that controls the spacing between line numbers and the text body).

Note: Bars can only be applied to complete paragraphs.
For bars between arbitrary points, use the \pkg{changebar} package.
Any limitations of the \pkg{lineno} package apply equally to the \pkg{vertbars} package.

\section*{Change History}

\begin{itemize}
\item[v1.0d 2026/10/21] New maintainer (\LaTeX{} Project)
\item[v1.0c 2018/12/12] Add \cs{par} before and after the \texttt{vertbar} environment to avoid confusing/unwanted behaviour with paragraphs
\item[v1.0b 2010/11/27] Fix behaviour with \texttt{switch} option
\item[v1.0a 2009/09/04] New maintainer (Will Robertson)
\end{itemize}



\section*{Licence and copyright}

This work may be modified and/or distributed under the terms and
conditions of the \LaTeX\ Project Public License\footnote{\url{https://www.latex-project.org/lppl.txt}}, version~1.3c or later (your choice).
The current maintainer of this work is the \LaTeX{} Project.

\bigskip
  \noindent
  Copyright Peter Wilson, 2000 \\
  Copyright Will Robertson, 2010--2018 \\
  Copyright \LaTeX{} Project, 2026

\end{document}
